\par\Needspace{6\baselineskip}\medskip
\hypertarget{native-085-article066}{}\pdfbookmark[2]{66.}{article-066}
\begin{center}66.\end{center}
Aus der Länge und Breite eines beliebigen Punktes auf der Himmelskugel erhält man seine Rektaszension und Deklination, indem man das sphärische Dreieck löst, das dieser Punkt mit den Nordpolen der Ekliptik und des Äquators bildet. Es seien $\epsilon$ die Schiefe der Ekliptik, $l$ die Länge, $b$ die Breite, $\alpha$ die Rektaszension und $\delta$ die Deklination. Die Seiten des Dreiecks sind dann $\epsilon$, $90^0-b$, $90^0-\delta$. Als Winkel gegenüber der zweiten und dritten Seite darf man $90^0+\alpha$, $90^0-l$ nehmen (sofern man den Begriff des sphärischen Dreiecks in seiner ganzen Allgemeinheit auffasst). Den dritten Winkel, der gegenüber der Seite $\epsilon$ liegt, setzen wir $=90^0-E$. Nach den Formeln aus Artikel 54 erhalten wir somit
\SourceLeaf{086}
\[
\begin{aligned}
\sin(45^0-\tfrac12\delta)\sin\tfrac12(E+\alpha)&=\sin(45^0+\tfrac12 l)\sin(45^0-\tfrac12(\epsilon+b))\\[.5em]
\sin(45^0-\tfrac12\delta)\cos\tfrac12(E+\alpha)&=\cos(45^0+\tfrac12 l)\cos(45^0-\tfrac12(\epsilon-b))\\[.5em]
\cos(45^0-\tfrac12\delta)\sin\tfrac12(E-\alpha)&=\cos(45^0+\tfrac12 l)\sin(45^0-\tfrac12(\epsilon-b))\\[.5em]
\cos(45^0-\tfrac12\delta)\cos\tfrac12(E-\alpha)&=\sin(45^0+\tfrac12 l)\cos(45^0-\tfrac12(\epsilon+b)).
\end{aligned}
\]
\noindent Die ersten beiden Gleichungen liefern $\tfrac12(E+\alpha)$ und $\sin(45^0-\tfrac12\delta)$. Die letzten beiden liefern $\tfrac12(E-\alpha)$ und $\cos(45^0-\tfrac12\delta)$. Aus $\tfrac12(E+\alpha)$ und $\tfrac12(E-\alpha)$ erhält man $\alpha$ und zugleich $E$. Aus $\sin(45^0-\tfrac12\delta)$ oder $\cos(45^0-\tfrac12\delta)$ bestimmt man $45^0-\tfrac12\delta$ und daraus $\delta$. Die Übereinstimmung beider Ergebnisse dient zur Kontrolle der Rechnung. Die Bestimmung der Winkel $\tfrac12(E+\alpha)$, $\tfrac12(E-\alpha)$ aus ihren Tangenten ist eindeutig, da sowohl der Sinus als auch der Kosinus des Winkels $45^0-\tfrac12\delta$ positiv ausfallen müssen.

Die differenziellen Änderungen der Größen $\alpha$, $\delta$ ergeben sich nach bekannten Grundsätzen aus den Änderungen von $l$, $b$ wie folgt.
\[
\begin{aligned}
\mathrm{d}\alpha&=\frac{\sin E\cos b}{\cos\delta}\,\mathrm{d}l-\frac{\cos E}{\cos\delta}\,\mathrm{d}b\\[.6em]
\mathrm{d}\delta&=\cos E\cos b\,\mathrm{d}l+\sin E\,\mathrm{d}b.
\end{aligned}
\]
